\documentclass[11pt]{article}
\usepackage[T1]{fontenc}
\usepackage{amsmath,amssymb,amsthm,mathtools,mathrsfs}
\usepackage[margin=24mm]{geometry}
\usepackage[colorlinks=true,linkcolor=blue,urlcolor=blue]{hyperref}
\newtheorem{theorem}{Theorem}[section]
\newtheorem{lemma}[theorem]{Lemma}
\newtheorem{corollary}[theorem]{Corollary}
\newcommand{\C}{\mathbb C}
\newcommand{\R}{\mathbb R}
\newcommand{\B}{\mathscr B}
\newcommand{\M}{\mathcal M}
\newcommand{\eps}{\varepsilon}
\newcommand{\Res}{\operatorname{Res}}
\newcommand{\im}{\operatorname{im}}
\newcommand{\Id}{\operatorname{Id}}
\title{The literal prime pulses observe every finite full zero packet}
\author{}
\date{23 September 2026}
\begin{document}
\maketitle

This note computes the exact map from the literal pulse family
FP4--6 of \cite{FP} to the original nonreduced zero packet of
\cite{PL}. It proves that an appropriate finite set of distinct
primes makes this map onto, without assuming that the bump factor
is a unit. Every local Taylor coefficient is retained. The final
sections construct both attained quotient metrics and identify the
whole complementary Weil contribution of their actual lifts.
The result concerns finite-packet observation; it is not positivity
of the original packet trace or of all Weil tests.

\section{The original packet, pulse, and Mellin conversion}

To distinguish the arithmetic completion from a pulse test, write
\(g_\zeta(s)=2\xi(s)\) for the function denoted \(g\) in \cite{PL}.
Thus
\[
 g_\zeta(s)=s(s-1)\pi^{-s/2}\Gamma(s/2)\zeta(s),
 \qquad g_\zeta(s)=g_\zeta(1-s)
               =\overline{g_\zeta(\bar s)}.
 \tag{PO1}
\]
Let \(Z\) be a finite nonempty set of actual nontrivial zeros,
with each full multiplicity \(m_\rho\) retained. Fix the original
nonzero leading coefficient \(c_h\) and define
\[
 \begin{gathered}
 h(s)=c_h\prod_{\rho\in Z}(s-\rho)^{m_\rho},\qquad
 d=\sum_{\rho\in Z}m_\rho,\qquad
 A_h=\C[s]/(h)\cong\prod_{\rho\in Z}\C[t_\rho]/(t_\rho^{m_\rho}),\\
 t_\rho=s-\rho,\qquad
 g_\zeta(\rho+t_\rho)=u_\rho(t_\rho)t_\rho^{m_\rho},
 \quad u_\rho(0)=2\lambda_\rho\ne0 .
 \end{gathered}
 \tag{PO2}\label{po:packet}
\]
The finite-dimensional coordinates throughout are the actual
Taylor coefficients of \(t_\rho^k\), \(0\le k<m_\rho\).
No unit \(u_\rho\), multiplicity, or leading coefficient is
replaced by a preferred value. For an entire function \(H\),
\[
 j_hH=\left(\sum_{k=0}^{m_\rho-1}
                    \frac{H^{(k)}(\rho)}{k!}t_\rho^k\right)_{\rho\in Z}.
 \tag{PO3}\label{po:jet}
\]
The human source for the full finite-jet spectral description is
Connes--Consani \cite{CC}, original equation \texttt{evalmaprho};
our coordinate factorials in PO3 are stated explicitly.

Fix a finite set \(I=\{p_1,\ldots,p_n\}\) of distinct primes,
where \(n\ge2\), and retain exactly the FP width and functions:
\[
 \begin{gathered}
 M_I=\prod_{i=1}^np_i,\qquad
 \eps_I=\frac1{16}\log(1+1/M_I),\qquad r_0=1,\quad r_i=p_i,\\
 b(v)=
 \begin{cases}e^{-1/(1-v^2)},&|v|<1,\\0,&|v|\ge1,\end{cases}
 \qquad
 w_\eps(t)=b(t/\eps),\quad
 z_\eps(t)=w_\eps''(t)-\frac14w_\eps(t),\\
 g_{I,c}(e^t)=\sum_{j=0}^n c_jz_{\eps_I}(t-\log r_j),
 \qquad
 N_I=\int_\R z_{\eps_I}(t)^2\,dt>0 .
 \end{gathered}
 \tag{PO4}\label{po:pulse}
\]
The smoothness, zero endpoint derivatives, and positivity of \(N_I\)
follow directly from the bump: all derivatives are an exponential
times a finite sum of rational powers of \(1-v^2\);
\(\int z_\eps w_\eps=-\int|w_\eps'|^2-\tfrac14\int w_\eps^2<0\).
The original logarithmic norm is
\[
 \|g_{I,c}\|_{L^2(dx/x)}^2=N_I\sum_{j=0}^n|c_j|^2 .
 \tag{PO5}\label{po:norm}
\]
Indeed distinct centres have distance greater than \(16\eps_I\):
for distinct positive integers \(a,b\),
\(|\log(a/b)|\ge\log(1+1/b)\), and every \(r_j<M_I\).
Their supports, each of radius \(\eps_I\), are disjoint.

The exact conversion to the original Mellin space is
\[
 F_{I,c}(x)=x^{-1/2}g_{I,c}(x),\qquad
 \M F(s)=\int_0^\infty F(x)x^s\,\frac{dx}{x}.
 \tag{PO6}\label{po:conversion}
\]
Every \(F_{I,c}\) is smooth with compact support in
\(\R_{>0}\), hence belongs to the original space
\[
 \B=\{F\in C^\infty(\R_{>0}):
       \sup_{x>0}|x^b(-x\partial_x)^kF(x)|<\infty
                \text{ for every }b\in\mathbb Z,\ k\ge0\}.
 \tag{PO7}
\]
This conversion does not preserve the unweighted \(L^2(dx/x)\)
norm of \(F\). Its exact isometry is
\(\int_0^\infty x|F_{I,c}(x)|^2\,dx/x=\|g_{I,c}\|^2\);
PO5 always refers to the original pulse norm.

Define the entire bump transform, with its actual amplitude,
\[
 W_\eps(z)=\int_\R w_\eps(t)e^{zt}\,dt
     =\eps\int_{-1}^1 b(v)e^{\eps zv}\,dv,\qquad
 a(s)=s(s-1),\qquad B_\eps(s)=a(s)W_\eps(s-\tfrac12).
 \tag{PO8}\label{po:bump}
\]
Two integrations by parts, with the literal zero endpoint terms,
give \(\int z_\eps(t)e^{zt}\,dt=(z^2-\tfrac14)W_\eps(z)\).
Translation by \(\log r_j\) multiplies this by \(r_j^z\).
Substitution \(x=e^t\) in PO6 therefore proves
\[
 \boxed{\quad
 \M F_{I,c}(s)=B_{\eps_I}(s)
             \sum_{j=0}^n c_j r_j^{\,s-1/2}.
 \quad}
 \tag{PO9}\label{po:transform}
\]
In particular \(\M F_{I,c}(0)=\M F_{I,c}(1)=0\).
In the original convention
\(\widehat g(\xi)=\int g(e^t)e^{-i\xi t}\,dt\) from
Connes \cite{Connes}, one has the exact signs
\[
 \M F(s)=\widehat g(i(s-\tfrac12)),\qquad
 \widehat g(-i(s-\tfrac12))=\M F(1-s).
 \tag{PO10}\label{po:Fourier}
\]
These two different arguments must not be interchanged without
the displayed reflection.

\section{All Taylor coefficients and the exact defect of a nonunit bump}

Put \(\alpha_\rho=\rho-\tfrac12\). For \(k\ge0\) define
\[
 \begin{aligned}
 w_{\rho,k}(\eps)
   &=\frac{W_\eps^{(k)}(\alpha_\rho)}{k!}
     =\frac{\eps^{k+1}}{k!}
        \int_{-1}^1v^k b(v)e^{\eps\alpha_\rho v}\,dv,\\
 b_{\rho,k}(\eps)
   &=\rho(\rho-1)w_{\rho,k}(\eps)
       +(2\rho-1)w_{\rho,k-1}(\eps)+w_{\rho,k-2}(\eps),
 \end{aligned}
 \tag{PO11}\label{po:taylor}
\]
where a coefficient with negative index is zero. These are the
full Taylor coefficients of \(W_\eps(s-\tfrac12)\) and
\(B_\eps(s)\) at the original point \(\rho\).

Let \(V_I:\C^{n+1}\to A_h\) have columns
\(v(r_j)=j_h(r_j^{s-1/2})\). In the actual coefficient basis,
let \(D_\eps\) be multiplication by \(j_hB_\eps\).
Its block at \(\rho\) is lower triangular, with
\((D_\eps)_{\rho;k,\ell}=b_{\rho,k-\ell}(\eps)\) for \(k\ge\ell\)
and zero otherwise. The full observation matrix is
\[
 \begin{gathered}
 E_I:\C^{n+1}\longrightarrow A_h,\qquad
 E_Ic=j_h\M F_{I,c},\qquad E_I=D_{\eps_I}V_I,\\
 (V_I)_{\rho,k;j}
       =r_j^{\alpha_\rho}\frac{(\log r_j)^k}{k!},\\
 (E_I)_{\rho,k;j}
       =r_j^{\alpha_\rho}\sum_{\ell=0}^k
           b_{\rho,k-\ell}(\eps_I)\frac{(\log r_j)^\ell}{\ell!}.
 \end{gathered}
 \tag{PO12}\label{po:matrix}
\]
Here \(\log r_0=0\); the zeroth power is one and positive powers
are zero. Formula PO12 retains every bump derivative, the quadratic
factor \(a\), all prime scalars, and every nilpotent power.

No invertibility has yet been assumed. Define
\[
 \nu_\rho(\eps)=
    \min\{m_\rho,\operatorname{ord}_{z=\alpha_\rho}W_\eps(z)\}.
 \tag{PO13}\label{po:vanishing}
\]
The function \(W_\eps\) is not identically zero, since
\(W_\eps(0)>0\), so its order at every point is finite.
Since \(a(\rho)\ne0\), multiplication by \(B_\eps\) on the
\(\rho\)-factor has the exact image, kernel, and quotient
\[
 \begin{gathered}
 \im D_{\eps,\rho}=(t_\rho^{\nu_\rho}),\qquad
 \ker D_{\eps,\rho}=(t_\rho^{m_\rho-\nu_\rho}),\\
 A_\rho/\im D_{\eps,\rho}\cong\C[t_\rho]/(t_\rho^{\nu_\rho}).
 \end{gathered}
 \tag{PO14}\label{po:nonunit}
\]
If \(\nu_\rho=0\), the displayed kernel is zero and quotient
zero; if \(\nu_\rho=m_\rho\), the image is zero and kernel all
of \(A_\rho\). To prove these statements, factor
\(B_\eps(\rho+t)=t^{\nu_\rho}u(t)\) with \(u(0)\ne0\)
when \(\nu_\rho<m_\rho\). Multiplication by \(u\) is invertible
by finite recursive division, and multiplication by \(t^{\nu_\rho}\)
has exactly the displayed coefficient kernel and image.
If \(\nu_\rho=m_\rho\), the multiplication is zero directly.

For an arbitrary finite prime set, write \(\mathcal V_I=\im V_I\).
The exact further coefficient and bump defects are
\[
 \begin{gathered}
 \im E_I=D_{\eps_I}\mathcal V_I,\\
 0\longrightarrow\ker V_I\longrightarrow\ker E_I
     \xrightarrow{\,V_I\,}\mathcal V_I\cap\ker D_{\eps_I}
       \longrightarrow0,\\
 0\longrightarrow D_{\eps_I}A_h/D_{\eps_I}\mathcal V_I
   \longrightarrow A_h/D_{\eps_I}\mathcal V_I
   \longrightarrow A_h/D_{\eps_I}A_h\longrightarrow0 .
 \end{gathered}
 \tag{PO15}\label{po:defects}
\]
Each exactness assertion follows by taking the indicated kernel
or image of its displayed linear map. If \(V_I\) is onto,
the last quotient has dimension \(\sum_\rho\nu_\rho(\eps_I)\);
thus PO14 is the precise missing-jet space of a nonunit bump.

\section{Shrinking the prescribed width controls the whole bump}

Write the unchanged bump moments
\[
 \mu_k=\int_{-1}^1v^k b(v)\,dv,\qquad
 \mu_0>0,\qquad \mu_{2k+1}=0,\qquad 0\le\mu_{2k}\le\mu_0 .
 \tag{PO16}
\]
Uniform convergence of the exponential series on \([-1,1]\)
gives
\[
 W_\eps(z)=\eps\sum_{k=0}^\infty
              \frac{\mu_{2k}(\eps z)^{2k}}{(2k)!},
 \qquad
 |W_\eps(z)-\eps\mu_0|
       \le\eps\mu_0(\cosh(\eps|z|)-1)
       \le\eps\mu_0(e^{\eps|z|}-1).
 \tag{PO17}\label{po:unit-control}
\]
The division by \(\eps\mu_0\) used to measure this estimate does
not change the pulse or any entry of PO12.
For all retained derivatives one also has
\[
 |w_{\rho,k}(\eps)|
   \le\frac{\eps^{k+1}\mu_0}{k!}e^{\eps|\alpha_\rho|}.
 \tag{PO18}\label{po:derivative-bound}
\]
This follows from PO11 and \(|v|\le1\).

Let \(R_Z=\max_{\rho\in Z}|\alpha_\rho|\).
If \(\eps R_Z<\log2\), then PO17 shows that
\[
 |W_\eps(\alpha_\rho)|
   \ge\eps\mu_0(2-e^{\eps R_Z})>0
                    \qquad(\rho\in Z).
 \tag{PO19}\label{po:unit-criterion}
\]
It follows that all \(\nu_\rho=0\) and \(D_\eps\) is invertible.
This proves invertibility; it does not replace the higher
coefficients by zero. Its complete determinant and inverse recursion are
\[
 \begin{gathered}
 \det D_\eps=\prod_{\rho\in Z}
       [\rho(\rho-1)W_\eps(\alpha_\rho)]^{m_\rho},\\
 d_{\rho,0}=b_{\rho,0}^{-1},\qquad
 d_{\rho,k}=-b_{\rho,0}^{-1}
                \sum_{j=1}^k b_{\rho,j}d_{\rho,k-j}
                   \quad(1\le k<m_\rho).
 \end{gathered}
 \tag{PO20}\label{po:inverse}
\]
The polynomial \(\sum d_{\rho,k}t_\rho^k\) is the exact inverse
of \(B_\eps\) in that factor, by multiplying the two finite series.

For the literal width of PO4,
\(\eps_I<1/(16M_I)\). Consequently the sufficient condition
\[
 M_I>R_Z/8
 \quad\Longrightarrow\quad
 \eps_I R_Z<1/2<\log2
 \tag{PO21}\label{po:width}
\]
proves the required unit property. The inequality
\(\log2=\int_0^1(1+t)^{-1}dt>1/2\) fixes the last constant.
Adding primes increases \(M_I\) without changing the defining
formula for the width. The actual new pulse has that smaller
width; it is not identified with the old pulse.

More precisely, for \(I\subseteq J\), let
\(\iota_{IJ}\) insert zero coefficients for the added primes.
The exponential maps satisfy \(V_J\iota_{IJ}=V_I\), and the
full jet maps satisfy
\[
 D_{\eps_I}E_J\iota_{IJ}=D_{\eps_J}E_I.
 \tag{PO22}\label{po:width-change}
\]
The multiplication operators commute, proving the equality.
If \(D_{\eps_I}\) is invertible, this is
\(E_J\iota_{IJ}=D_{\eps_J}D_{\eps_I}^{-1}E_I\).
Thus PO22 is the exact comparison when the bump width changes;
all unit Taylor data in the comparison are given by PO11 and PO20.

\section{Prime logarithms determine every exponential polynomial}

We give the uniqueness argument in full. It uses neither a
prime-number theorem nor a spacing assumption on the primes.
Its circle identity is the classical Jensen formula
\cite{Jensen}; its arithmetic input is Euler's reciprocal-prime
argument \cite{Euler}. Both are derived completely below.
For these historical works the cited archive metadata, rather
than an unavailable original TeX source, was consulted.

\begin{lemma}[An exponential polynomial has at most linearly many
zeros in growing discs]\label{lem:zero-count}
If
\[
 f(z)=\sum_{\ell=1}^r P_\ell(z)e^{\alpha_\ell z}
 \tag{PO23}\label{po:exp-polynomial}
\]
is a nonzero entire function, with polynomial \(P_\ell\), then
there is a constant \(C_f\) such that its number of zeros in
\(|z|\le T\), counted with multiplicity, is at most
\(C_f(T+1)\) for \(T\ge1\).
\end{lemma}
\begin{proof}
With \(k=\max\deg P_\ell\) and \(A=\max|\alpha_\ell|\), one has
\(|f(z)|\le C(1+|z|)^k e^{A|z|}\).
Factor its possible zero at the origin as \(f(z)=z^q f_0(z)\),
where \(f_0(0)\ne0\).
Choose \(r\) with \(2T<r<3T\) whose circle contains no zero.
Such radii exist because the zeros of a nonzero entire function
are isolated. List the finitely many zeros \(\zeta_j\) of \(f_0\)
inside that circle, with multiplicities. After dividing out
\(\prod_j(z-\zeta_j)\), the remaining holomorphic function has
no zeros on the closed disc. Its logarithmic modulus is harmonic
and satisfies the circle mean-value identity. For \(|\zeta_j|<r\),
the average of \(\log|re^{i\theta}-\zeta_j|\) is \(\log r\):
expand \(\log(1-\zeta_je^{-i\theta}/r)\) into its uniformly
convergent power series and average its real part.
Subtracting the central value therefore gives the exact formula
\[
 \sum_{|\zeta_j|<r}\log\frac r{|\zeta_j|}
  =\frac1{2\pi}\int_0^{2\pi}
       \log|f_0(re^{i\theta})|\,d\theta-\log|f_0(0)|.
 \tag{PO24}\label{po:circle-count}
\]
Each zero with \(|\zeta_j|\le T\) contributes at least \(\log2\).
The right side is at most
\(\log C+k\log(1+3T)+3AT-\log|f_0(0)|\);
the factor \(r^{-q}\) can only decrease this bound for \(T\ge1\).
This is at most a constant times \(T+1\).
Adding the \(q\) zeros at zero proves the lemma.
\end{proof}

\begin{lemma}[Prime reciprocals diverge]\label{lem:Euler}
For every finite set \(S\) of primes,
\(\sum_{p\notin S}p^{-1}=\infty\).
\end{lemma}
\begin{proof}
It suffices to prove divergence over all primes.
If that sum were finite, then
\[
 \log\prod_{p\le X}(1-p^{-1})^{-1}
  =\sum_{p\le X}\int_0^{1/p}\frac{dt}{1-t}
  \le2\sum_p p^{-1}
 \tag{PO25}
\]
would bound every finite product independently of \(X\).
Expanding its finitely many geometric factors, with nonnegative
terms and unique factorization, shows that this product contains
every term \(1/m\) for \(1\le m\le X\), since each prime factor
of such an \(m\) is at most \(X\).
Thus it is at least \(\sum_{m\le X}1/m\), which tends to
infinity by comparison with \(\int_1^X dt/t\).
This contradiction proves divergence. Removing finitely many
terms leaves divergence unchanged.
\end{proof}

\begin{theorem}[Uniqueness on every cofinite prime-log set]
\label{thm:prime-uniqueness}
Let \(\alpha_\ell\) be pairwise distinct complex numbers,
and let \(P_\ell\) be polynomials. If the function PO23
vanishes at \(\log p\) for all primes outside a finite set,
then every \(P_\ell\) is zero.
\end{theorem}
\begin{proof}
First suppose \(f\) were nonzero. Lemma \ref{lem:zero-count}
would give at most \(C(k+2)\) distinct prime logarithms in
each interval \([k,k+1)\), since those are zeros in the disc
of radius \(k+1\).
It would follow that
\[
 \sum_{p\notin S}\frac1p
   =\sum_{p\notin S}e^{-\log p}
   \le C\sum_{k=0}^\infty(k+2)e^{-k}<\infty,
 \tag{PO26}
\]
contradicting Lemma \ref{lem:Euler}. Hence \(f\equiv0\).

To exclude cancellation between distinct exponents, suppose
\(P_{\ell_0}\ne0\) and set
\(m_\ell=\deg P_\ell+1\) for each nonzero polynomial.
Apply
\(\prod_{\ell\ne\ell_0}(\partial_z-\alpha_\ell)^{m_\ell}\)
to \(f\equiv0\). Every term except the \(\ell_0\) term is
annihilated. The latter becomes \(e^{\alpha_{\ell_0}z}\)
times
\(\prod_{\ell\ne\ell_0}
  (\partial_z+\alpha_{\ell_0}-\alpha_\ell)^{m_\ell}
    P_{\ell_0}(z)\).
On the finite-dimensional polynomial space of degree at most
\(\deg P_{\ell_0}\), each factor has a nonzero constant
diagonal and a strictly degree-lowering remainder. It is
invertible by a finite geometric inverse. The resulting
polynomial is therefore nonzero, a contradiction.
Thus all \(P_\ell\) vanish.
\end{proof}

\section{A finite literal pulse family is onto all original jets}

\begin{theorem}[Full-jet surjectivity with the original width]
\label{thm:onto}
For every finite packet PO2, there is a finite set \(I\) of
distinct primes, \(|I|\ge2\), such that the literal observation
\(E_I\) in PO12 is onto \(A_h\).
The set can be required to contain any prescribed finite prime
set. It can also be chosen with
\(\eps_I R_Z<\log2\), so its bump factor is a proved unit
with all its coefficients still given by PO11.
\end{theorem}
\begin{proof}
For every prime \(p\), consider the full vector
\[
 v(p)=
 \left(p^{\alpha_\rho}
            \frac{(\log p)^k}{k!}\right)_{\rho\in Z,\ 0\le k<m_\rho}.
 \tag{PO27}\label{po:prime-columns}
\]
These vectors span \(\C^d\), even after any finite set of
primes is removed. If their span were proper, elementary
finite-dimensional linear algebra would give a nonzero
complex linear functional annihilating it. In coordinates,
its prime values would be
\[
 \sum_{\rho\in Z}p^{\alpha_\rho}
       \sum_{k=0}^{m_\rho-1}\ell_{\rho,k}
                           \frac{(\log p)^k}{k!}=0 .
 \tag{PO28}
\]
The exponents \(\alpha_\rho\) are distinct. Theorem
\ref{thm:prime-uniqueness}, applied to \(z=\log p\), forces
every coefficient \(\ell_{\rho,k}\) to vanish, a contradiction.

Select \(d\) distinct prime vectors forming a basis. This
selection can be performed successively: whenever the
chosen span is proper, the spanning assertion supplies
another prime outside it, also outside any prescribed finite
exclusion. Take the union of these primes with the prescribed
finite set. If needed, add primes so that the total cardinality
is at least two and its product exceeds \(R_Z/8\). There
are arbitrarily large primes by Lemma \ref{lem:Euler}, so
finitely many additions suffice.

Call the resulting set \(I\). The columns of \(V_I\) still
contain the selected invertible \(d\)-by-\(d\) submatrix;
enlarging \(I\) does not alter these exponential columns.
Its width now satisfies PO21, so \(D_{\eps_I}\) is
invertible by PO19. Therefore \(E_I=D_{\eps_I}V_I\) is
onto. Its retained invertible minor has determinant
\[
 \left(\prod_{\rho\in Z}
 [\rho(\rho-1)W_{\eps_I}(\alpha_\rho)]^{m_\rho}\right)
       \det[v(p_1')\ \cdots\ v(p_d')]\ne0,
 \tag{PO29}\label{po:minor}
\]
where \(p_1',\ldots,p_d'\) are the selected basis primes.
This proves surjectivity with the literal new width and every
local factor included.
\end{proof}

The proof provides existence of a finite spanning minor and an
explicit width threshold. It does not assert a uniform bound for
the size of those basis primes. The observation does not require
a constant-column contribution: the \(d\) selected prime columns
already form a basis; the original \(r_0=1\) column remains present.

\subsection{The actual cohomological map and prime convention}

Retain the original source
\(V=\{\phi\in\mathcal S(\R):\phi\text{ even},
\ \phi(0)=0,\ \int\phi=0\}\),
\(\Theta\phi(x)=2\sum_{n\ge1}\phi(nx)\), and \(Q=\B/\Theta V\).
For \(\Re s>1\), substitution gives
\(\M\Theta\phi(s)=2\zeta(s)M_+\phi(s)\), with
\(M_+\phi(s)=\int_0^\infty\phi(x)x^s\,dx/x\).
The latter is holomorphic for \(\Re s>-2\):
evenness and \(\phi(0)=0\) imply \(\phi(x)=O(x^2)\) at zero,
while the source is rapidly decreasing at infinity.
The identity therefore continues to the critical strip and
shows that all full actual zero jets of \(\Theta\phi\) vanish.
Consequently the exact observation factors as
\[
 \C^{n+1}\xrightarrow{\,c\mapsto F_{I,c}\,}
 \B_{\rm pr}\longrightarrow Q
       \xrightarrow{\,\bar j_h\,}A_h,\qquad
 \bar j_h[F]=j_h\M F,\qquad
 \bar j_h[F_{I,c}]=E_Ic .
 \tag{PO30}\label{po:Qmap}
\]
Here \(\B_{\rm pr}\) imposes the two Mellin endpoint zeros,
already verified in PO9. Surjectivity of the composite does
not identify its kernel with the theta relation space:
\(\ker E_I\) means exactly vanishing at this finite packet.

The full original prime action on \(F\) is
\(U_pF(x)=F(x/p)\), with \(\M U_pF(s)=p^s\M F(s)\).
On a literal pulse translation \(g(x)\mapsto g(x/p)\),
conversion PO6 gives instead
\[
 x^{-1/2}g(x/p)=p^{-1/2}U_pF(x),\qquad
 \bar j_h[p^{-1/2}U_pF]
    =p^{-1/2}\Pi_p\,\bar j_h[F],\quad
 \Pi_p=M_{p^s}.
 \tag{PO31}\label{po:prime-factor}
\]
The local matrix is
\(p^{\rho-1/2}\sum_{k<m_\rho}
(\log p)^kN_\rho^k/k!\), where \(N_\rho\) multiplies by
\(t_\rho\). This is the exact factor connecting the two
original actions.
PO30--31 are assertions on the original arithmetic coefficient
spaces. Their finite observation has the following explicit
support diagram, rather than a claim that every support stalk
has constant coefficients. Let \(\mathcal J=\{0,\ldots,n\}\)
be its source indices and
\(\mathcal K=\{(\rho,k):\rho\in Z,\ 0\le k<m_\rho\}\)
its target indices. Retain a declared incidence set
\(\mathcal D\subseteq\mathcal K\times\mathcal J\) containing
every nonzero entry of the actual matrix \(E_I\); any additional
declared zero entries remain in \(\mathcal D\).
The exact support trajectories and coefficient maps are
\[
 \begin{gathered}
 \mathcal L_{\mathcal D}
  =\{(B,C):B\subseteq\mathcal J,\ C\subseteq\mathcal K,\
                         \mathcal D(B)\subseteq C\},\\
 E_{B,C}:\C^B\longrightarrow\C^C,\qquad
 E_{B,C}=(E_I)_{C,B},\\
 j_{CC'}E_{B,C}=E_{B',C'}j_{BB'}
       \quad\text{for }B\subseteq B',\ C\subseteq C'.
 \end{gathered}
 \tag{PO54}\label{po:support-square}
\]
Here \(j\) means extension by zero. To prove the last square,
apply it to a coordinate vector at \(j\in B\).
Every nonzero row coefficient \(E_{\alpha j}\) has
\(\alpha\in\mathcal D(B)\subseteq C\); hence the two sides
have exactly the same retained coefficients in \(C'\).
Declared zero entries change neither equality nor their
support labels. Componentwise unions and intersections preserve
the trajectory condition, so this is the original finite
support-trajectory construction. Its top is
\((\mathcal J,\mathcal K)\), where the coefficient map is
exactly \(E_I:\C^{n+1}\to A_h\).
Each smaller coefficient space is the indicated coordinate
space, and need not have the dimension of its top.
Thus PO12 is the terminal arithmetic fibre of this explicit
diagram; it has not replaced the other support stalks.
This support indexing is distinct from the prime-face index
\(A\subseteq I\) used for the Weil matrices below.

\section{Two attained positive metrics, with their exact sections}

Fix a prime set from Theorem \ref{thm:onto}, and write
\(E=E_I\), \(N=N_I\). All adjoints in this section use the
explicit coefficient bases \(c_0,\ldots,c_n\) and
\(t_\rho^k\). The standard coordinate inner product is used
only to write these matrices; the source pulse norm remains
the literal \(N\Id\) from PO5.

For completeness retain the exact Hermitian FP form.
Let \(C_\eps(t)=\int z_\eps(u)z_\eps(u-t)\,du\),
\(K(t)=e^{-t/2}/(1-e^{-2t})\),
\(c_{\rm arch}=\log(4\pi)+\gamma_E\), and
\[
 \begin{aligned}
 A_\eps&=c_{\rm arch}N+
       2\int_0^\infty[C_\eps(t)-Ne^{-t/2}]K(t)\,dt,\\
 J_\eps(a)&=\int_{a-2\eps}^{a+2\eps}C_\eps(t-a)K(t)\,dt,\\
 (\mathsf H_\eps)_{ij}
   &=\begin{cases}J_\eps(|\log r_i-\log r_j|),&i\ne j,\\0,&i=j.\end{cases}
 \end{aligned}
 \tag{PO32}
\]
For each retained prime face \(A\subseteq I\), the full matrix is
\[
 \mathsf W_A=-A_\eps\Id-\mathsf H_\eps
    -N\sum_{p_i\in A}\frac{\log p_i}{\sqrt{p_i}}
                         (\mathsf e_{0i}+\mathsf e_{i0}).
 \tag{PO33}\label{po:FPmatrix}
\]
In particular \(\mathsf W=\mathsf W_I\) is the actual complete
Weil matrix of these pulse functions. FP18--43 of \cite{FP}
prove, including the entire Archimedean subtraction,
\[
 Q(g_{I,c})=c^*\mathsf Wc,\qquad
 \mathsf W_A\ge\mu_*N\Id,\qquad
 \mu_*=\frac{6600077509}{144027072000}>\frac1{22}.
 \tag{PO34}\label{po:coercive}
\]
For every nonzero \(c\), the inequality proved there is strict.
The complete companion source supplies that proof; it uses no
positivity assumption on a zero packet.
Thus \(\mathsf W_A\) is positive definite on the whole
coefficient space for every face, without suppressing any
cross coefficient in PO33.

\begin{theorem}[The two minimum sections]\label{thm:minima}
The minimum section for the original pulse norm is
\[
 L_2=E^*(EE^*)^{-1},\qquad
 EL_2=\Id,\qquad
 G_2=N(EE^*)^{-1}.
 \tag{PO35}\label{po:L2}
\]
For every \(u\in A_h\),
\(\min_{Ec=u}N\|c\|_2^2=u^*G_2u\), attained uniquely at \(L_2u\).
Its pulled-back full Weil form is
\(G_{\rm lift}=L_2^*\mathsf W L_2\), which is positive definite.

The minimum section for the full Weil coefficient form is instead
\[
 G_W=(E\mathsf W^{-1}E^*)^{-1},\qquad
 L_W=\mathsf W^{-1}E^*G_W,\qquad
 EL_W=\Id,\qquad L_W^*\mathsf W L_W=G_W .
 \tag{PO36}\label{po:LW}
\]
It gives
\[
 \min_{Ec=u}c^*\mathsf Wc=u^*G_Wu,\qquad
 G_{\rm lift}\ge G_W\ge\mu_*N(EE^*)^{-1}.
 \tag{PO37}\label{po:metrics}
\]
The same constructions hold for each face after replacing
\(\mathsf W\) by its exact matrix \(\mathsf W_A\).
\end{theorem}
\begin{proof}
Since \(E\) is onto, \(EE^*\) is positive definite:
\(u^*EE^*u=\|E^*u\|^2=0\) implies that \(u\) annihilates
the image of \(E\), hence \(u=0\).
Thus PO35 is defined and \(EL_2=\Id\).
Every solution is \(c=L_2u+k\), \(k\in\ker E\);
\(\langle L_2u,k\rangle=u^*(EE^*)^{-1}Ek=0\).
Expanding the norm proves the asserted minimum and its uniqueness.
It also gives \(L_2^*L_2=(EE^*)^{-1}\).

Positive definiteness of \(\mathsf W\) follows from PO34.
Therefore \(E\mathsf W^{-1}E^*\) is positive definite by
the same argument, and PO36 is defined.
For \(k\in\ker E\),
\((L_Wu)^*\mathsf Wk=u^*G_WEk=0\).
Every solution now satisfies the exact decomposition
\[
 c^*\mathsf Wc
   =u^*G_Wu+k^*\mathsf Wk,\qquad c=L_Wu+k,\quad Ek=0.
 \tag{PO38}\label{po:minimum-decomp}
\]
This proves the unique attained minimum and PO36.
Using \(L_2u\) as a competing solution proves
\(G_{\rm lift}\ge G_W\). For every solution \(Ec=u\),
PO34 and PO35 give
\(c^*\mathsf Wc\ge\mu_*N\|c\|^2
\ge\mu_*N u^*(EE^*)^{-1}u\).
Taking the attained minimum proves PO37.
For \(u\ne0\), \(L_Wu\ne0\), so the last inequality is
strict by the strict version of PO34.
Replacing \(\mathsf W\) by \(\mathsf W_A\) uses precisely
the same proof.
\end{proof}

These are different specified metrics. Their connecting identity is
\[
 G_{\rm lift}-G_W
  =(L_2-L_W)^*\mathsf W(L_2-L_W),\qquad
 E(L_2-L_W)=0 .
 \tag{PO39}\label{po:metric-difference}
\]
It follows by applying the orthogonality in PO38 to
\(L_2u=L_Wu+(L_2-L_W)u\) and polarizing.
Thus neither minimum section is silently substituted for the other.

\section{The exact complementary Weil contribution}

Now require \(Z\) to be stable under
\(\iota\rho=1-\bar\rho\), as in the original reflected packet.
For germs \(u,v\in A_h\), retain the original raw residue form
and its Jacobian contraction:
\[
 \begin{gathered}
 u^\dagger(s)=\overline{u(1-\bar s)},\qquad
 R_Z(u,v)=\sum_{\rho\in Z}
       \Res_\rho\frac{u^\dagger(s)v(s)}{g_\zeta(s)}\,ds,\qquad
 J_\zeta=M_{g_\zeta'},\\
 T_Z(u,v)=R_Z(u,J_\zeta v)
       =\sum_{\rho\in Z}m_\rho
                  \overline{u(\iota\rho)}v(\rho).
 \end{gathered}
 \tag{PO40}\label{po:trace}
\]
The equality follows from
\(g_\zeta'/g_\zeta=m_\rho/t_\rho+u_\rho'/u_\rho\):
only the simple pole contributes.
The full residue block, with row \(t_{\iota\rho}^j\) and
column \(t_\rho^k\), is
\[
 (-1)^j[t_\rho^{m_\rho-1-j-k}]\,u_\rho(t_\rho)^{-1}.
 \tag{PO41}\label{po:raw-matrix}
\]
Its antidiagonal is \((-1)^j/(2\lambda_\rho)\ne0\);
hence the raw pairing remains perfect. The Hermitian matrix
of \(T_Z\), denoted \(\mathsf W_Z\), has entries
\[
 (\mathsf W_Z)_{\sigma,k;\rho,\ell}
     =m_\rho\,\mathbf1_{\sigma=\iota\rho}
                \mathbf1_{k=0}\mathbf1_{\ell=0}.
 \tag{PO42}\label{po:WZ}
\]
All higher-jet coordinates remain in the packet even though
this particular contraction vanishes on them. The multiplicities
on a reflected pair agree by PO1.

We use the exact primitive packet section of \cite{PL}, and
recall its full construction to specify the complementary test.
Let \(U_h=g_\zeta/h\), let \(u_h=j_hU_h\), an actual unit of
\(A_h\), and set
\[
 R_u=\operatorname{rem}_h([a]^{-1}u_h^{-1}u),\qquad
 \phi_*(x)=(4\pi^2x^4-6\pi x^2)e^{-\pi x^2},\qquad D=-x\partial_x .
 \tag{PO43}
\]
The remainder has degree less than \(d\); both inverse units
are taken with their full Taylor expansions.
For \(H\in\B\) with \(\M H(\rho)=0\), put
\[
 S_\rho H(x)=x^{-\rho}\int_x^\infty H(t)t^{\rho-1}\,dt,
 \qquad S_h=c_h^{-1}\prod_{\rho\in Z}S_\rho^{m_\rho},\qquad
 \operatorname{Prim}_h(u)=a(D)S_h\Theta(R_u(D)\phi_*).
 \tag{PO44}\label{po:Prim}
\]
These expressions are well defined on the displayed arguments.
For clarity, differentiating the integral gives
\((D-\rho)S_\rho H=H\).
The zero Mellin moment lets one replace the tail near zero by
minus the integral from zero to \(x\). Rapid decrease of \(H\)
at both ends then bounds each weighted Euler derivative of
\(S_\rho H\); so \(S_\rho H\in\B\).
Its Mellin transform is \(\M H(s)/(s-\rho)\), by integration
by parts. Repeating removes the full zero orders, and uniqueness
of the rapidly decreasing solution proves that the factors commute.
Gaussian integration gives
\(\M\Theta\phi_*=g_\zeta\); multiplying by \(R_u(s)\) supplies
all required zeros of \(h\).
It follows that
\[
 \M\operatorname{Prim}_h(u)(s)=a(s)U_h(s)R_u(s),\qquad
 j_h\M\operatorname{Prim}_h(u)=u .
 \tag{PO45}\label{po:Prim-properties}
\]
The section is primitive at \(0,1\), and its transform has
the full actual zero multiplicity at every zero outside \(Z\).
This verifies the exact properties used below, including the
original factor \(c_h\), units \(u_h\), and coefficients of \(\phi_*\).

For \(F,H\in\B\), retain the whole reflected form
\[
 \mathcal W(F,H)=\sum_{\rho\text{ actual}}
     m_\rho\,\overline{\M F(1-\bar\rho)}\,\M H(\rho).
 \tag{PO46}\label{po:global}
\]
The series converges absolutely: integration by parts in
\(\log x\) gives arbitrary inverse powers of \(|\Im\rho|\)
uniformly in the critical strip. The zero-count bound
\(N(T)=O(T\log T)\) needed for summation is proved directly
from the unchanged Gaussian seed in \cite{PL}, PL50.
For the pulse conversion PO6, the original explicit formula
of \cite{Connes} gives
\[
 \mathcal W(F_{I,c},F_{I,d})=c^*\mathsf Wd.
 \tag{PO47}\label{po:Weil-conversion}
\]
To check the sign conversion, its Fourier zero argument is
\((\rho-\tfrac12)/i=-i(\rho-\tfrac12)\).
For a convolution square, the value there is
\(\M F(1-\rho)\overline{\M F(\bar\rho)}\), by PO10 and
the conjugate Fourier product. Reindex the entire zero
multiset by \(\rho\mapsto1-\rho\), which preserves multiplicity
by PO1; the result is exactly PO46.
Both pole terms vanish by PO9. The sesquilinear identity follows
by polarization, with conjugate linearity in the first slot.
Thus PO47 invokes the actual human explicit-formula convention,
not a new choice of trace sign.

Define the original residual map on all pulse coefficients by
\[
 \mathcal R_I(c)=F_{I,c}-\operatorname{Prim}_h(Ec)\in\B_{\rm pr}.
 \tag{PO48}\label{po:residual}
\]
It has every full jet zero on \(Z\), by PO12 and PO45.
At every actual zero outside \(Z\), the packet section vanishes
to full order, so the residual retains exactly the original
pulse's jets there. This is an actual function-space map,
not a formal difference of assigned quadratic forms.

\begin{theorem}[Exact paired residual and attained packet metrics]
\label{thm:residual}
For every coefficient pair \(c,d\),
\[
 \begin{gathered}
 \mathcal W(\operatorname{Prim}_h(Ec),\mathcal R_I(d))
 =\mathcal W(\mathcal R_I(c),\operatorname{Prim}_h(Ed))=0,\\
 \mathcal W(\mathcal R_I(c),\mathcal R_I(d))
      =c^*(\mathsf W-E^*\mathsf W_ZE)d .
 \end{gathered}
 \tag{PO49}\label{po:orthogonal}
\]
For any right inverse \(L\) of \(E\), put
\(\mathcal R_L(u)=F_{I,Lu}-\operatorname{Prim}_h(u)\).
Its exact whole complementary matrix is
\[
 \mathcal W(\mathcal R_L(u),\mathcal R_L(v))
       =u^*\mathsf D_Lv,\qquad
 \boxed{\mathsf D_L=L^*\mathsf W L-\mathsf W_Z.}
 \tag{PO50}\label{po:complement}
\]
In particular
\(\mathsf D_{L_2}=G_{\rm lift}-\mathsf W_Z\) and
\(\mathsf D_{L_W}=G_W-\mathsf W_Z\), with
\[
 \mathsf D_{L_2}\ge\mathsf D_{L_W}
       \ge\mu_*N(EE^*)^{-1}-\mathsf W_Z .
 \tag{PO51}\label{po:residual-lower}
\]
\end{theorem}
\begin{proof}
At a zero in \(Z\), every value of the residual vanishes,
including at its reflected partner because \(Z\) is stable.
At any zero outside \(Z\), every value of the primitive section
vanishes, also at its reflected partner.
Each summand of either cross form is therefore zero individually.
The pure section term of PO46 reduces exactly to \(T_Z(Ec,Ed)\)
by PO45, with no omitted outside contribution.
Expanding \(F_{I,c}=\operatorname{Prim}_h(Ec)+\mathcal R_I(c)\)
in PO47 proves PO49. Substitution \(c=Lu,d=Lv\) proves PO50.
The matrix inequalities are precisely PO37 with
\(\mathsf W_Z\) subtracted; hence PO51 follows.
\end{proof}

The residual class \([\mathcal R_L(u)]\) lies in
\(\ker(\bar j_h:Q\to A_h)\) by PO30.
It is not asserted to be a theta relation. Thus PO48 is also
the exact cohomological difference between a pulse lift and
the same-jet primitive packet section.

\begin{corollary}[Numerical lower bound in a negative packet direction]
\label{cor:negative}
For every nonzero \(u\) with \(u^*\mathsf W_Zu<0\),
the minimum-Weil lift satisfies
\[
 \mathcal W(\mathcal R_{L_W}(u),\mathcal R_{L_W}(u))
 >\mu_*N\,u^*(EE^*)^{-1}u-u^*\mathsf W_Zu>0.
 \tag{PO52}\label{po:negative}
\]
The same lower bound holds for the minimum-pulse-norm lift \(L_2\).
\end{corollary}
\begin{proof}
Use the strict version of PO37 for \(u\ne0\) in PO50,
then the fact that \((EE^*)^{-1}\) is positive definite.
For \(L_2\), use \(\mathsf D_{L_2}\ge\mathsf D_{L_W}\).
\end{proof}

More explicitly, on any nonfixed reflected pair
\(\{\rho,\iota\rho\}\) of common multiplicity \(m\), retain the
full coordinate vector \(u=\eta\), whose two constant
coefficients are \(1,-1\) and whose other coordinates are zero.
The original trace block is
\(m\left(\begin{smallmatrix}0&1\\1&0\end{smallmatrix}\right)\);
hence \(\eta^*\mathsf W_Z\eta=-2m\).
The exact numerical lower bound for its complementary
whole-zero contribution is
\[
 \mathcal W(\mathcal R_{L_W}(\eta),\mathcal R_{L_W}(\eta))
 >2m+
    \frac{6600077509}{144027072000}
            N\,\eta^*(EE^*)^{-1}\eta .
 \tag{PO53}\label{po:pair-bound}
\]
All entries of \(E\) in this bound are the full bump-and-prime
Taylor factors of PO11--12; \(N\) is the unchanged integral PO4.
For a scalar multiple \(a\eta\), both sides scale by \(|a|^2\).
This computes the exact compensation for an indefinite packet
direction whenever that direction occurs in the original packet;
it does not assert existence of an off-critical zero.
For a nonzero higher-jet vector with all constant coefficients
zero, PO42 instead gives \(T_Z(u,u)=0\), and the same proof gives
a strictly positive complementary lower bound
\(\mu_*N u^*(EE^*)^{-1}u\).
The raw form PO41 continues to detect those jets.

The onto map and the positive pulse form therefore coexist
through the explicit residual PO48. Neither surjectivity nor
the attained positive metric \(G_W\) identifies it with
the original Jacobian trace \(\mathsf W_Z\).
Their exact difference is the matrix and full zero sum PO50.

\begin{thebibliography}{9}
\bibitem{Jensen}
Johan Ludwig William Valdemar Jensen,
\emph{Sur un nouvel et important th\'eor\`eme de la th\'eorie des
fonctions}, Acta Mathematica \textbf{22} (1899), 359--364,
\href{https://doi.org/10.1007/BF02417878}{doi:10.1007/BF02417878}.
The \href{https://archive.ymsc.tsinghua.edu.cn/pacm_journal?journal=Acta+Mathematica&volume=22&year=1899}{journal archive metadata}
was checked; no original TeX or PDF was read.
PO24 is derived in full here from the circle mean-value identity.
\bibitem{Euler}
Leonhard Euler, \emph{Variae observationes circa series infinitas},
Enestr\"om E72, Commentarii academiae scientiarum Petropolitanae
\textbf{9} (1744), 160--188; written 1737.
The \href{https://scholarlycommons.pacific.edu/euler-works/72/}{Euler archive record}
was checked for attribution and metadata; no original TeX or
PDF was read. The reciprocal-prime divergence needed here
is proved completely in PO25.
\bibitem{Connes}
Alain Connes, \emph{The Riemann Hypothesis: Past, Present and a
Letter Through Time},
\href{https://arxiv.org/abs/2602.04022v1}{arXiv:2602.04022v1},
original author source \texttt{rhready.tex}, section
\texttt{sectweilpos}, lines 850--864, equations \texttt{wp}
and \texttt{arch}. The positivity criterion is attributed
there to Andr\'e Weil. PO10 and PO47 retain its Fourier and
Haar conventions explicitly.
\bibitem{CC}
Alain Connes and Caterina Consani,
\emph{Hochschild homology, trace map and zeta cycles},
\href{https://arxiv.org/abs/2207.10419}{arXiv:2207.10419},
original author source \texttt{Hochschild-zeta.tex},
lines 458--464, equation \texttt{evalmaprho} in the proof of
Proposition \texttt{laplacian}; full multiplicity jets and the
triangular action. PO2--3 state all coordinates used here.
\bibitem{FP}
Companion derivation, \emph{The complete Hermitian pulse matrix
and a uniform finite-prime Weil bound},
\texttt{FINITE\_PRIME\_PULSE\_BOUND.tex}, complete source
retained alongside this note, FP4--6, FP18--28 and FP29--43.
It proves the literal norm, all face matrices and the exact
constant used in PO34.
Its original two-prime programme input is the
\href{https://github.com/KokunoYumeto/zeta-function-research-reader/blob/3c78cfbd9e194d35117b33277c421c6564ce3dbf/workbenches/splitzero-tandem/continuations/20260923-full-prime-support/033/MIXED_PRIME_ENERGY_AND_BOUNDARY.tex}{complete pinned edition 033},
FWP1--24 and FWJ1--8.
\bibitem{PL}
Companion derivation, \emph{Primitive representatives of the
complete arithmetic jet packets},
\texttt{PRIMITIVE\_PACKET\_LIFT.tex}, complete source retained
alongside this note, PL10--18, PL23--28 and PL50.
PO43--45 repeat the full packet section and its relevant proof.
Its original programme convention is the
\href{https://github.com/KokunoYumeto/zeta-function-research-reader/blob/4ba9285b66af58a6d58fc502a494c1fb45ca5b79/workbenches/tau-base-cohomology/TAU_BASE_MODEL.md}{pinned tau-base source},
Sections 2--6.
\end{thebibliography}
\end{document}
